\section{Introduction}
\label{sec:companion-introduction}

How much can economic growth increase if AI systems help develop better AI?
Answering this question requires connecting the productivity of AI research
to the allocation of resources in the rest of the economy. Research competes
with current consumption, physical investment, and the compute used to supply
AI services. Its scale depends on the profits that improvements can earn and
on the interest rate at which those profits are discounted.

I study these decisions in an endogenous-growth model with a representative
household, competitive final-good firms, and a monopolistic AI developer.
Final production combines capital, labor, and AI services. The developer
uses AI research services to raise the efficiency with which compute produces
both production and research services. Research is autonomous: the model
describes an economy in which human researchers are no longer a necessary
current input into AI development. It does not describe how that transition
to autonomous research occurs.

This paper develops the uncapped unit-elastic benchmark in Appendix C of
\citet{pardo2026future}. That paper compares substitution regimes under a
finite AI-efficiency frontier. Here I fix the elasticity of substitution at
one and remove that frontier. I derive an equilibrium with constant growth
rates and endogenous allocation shares, rather than assume fixed saving and
research expenditure. This restriction isolates a tractable quantitative
question; it is not a claim that aggregate AI--labor substitution has been
estimated to equal one.

I distinguish inherited macroeconomic reference values from uncertain AI
parameters. The analytical results identify which parameters affect permanent
growth and which affect reference levels. This first draft establishes that
benchmark and the calibration framework. Empirically disciplined ranges for
the AI parameters and quantitative sensitivity exercises remain to be developed.
The calculations do not assume that the present world economy lies on the
model's balanced-growth path.
